\subsection{Возможные направления дальнейшего развития решения}

\subsubsection{Поддержка краевых случаев}

Компенсация ряда часто встречающихся краевых случаев может быть оптимизирована с помощью дополнительного анализа или добавления специализированных записей, компенсация которых может быть упрощена. Примеры возможных оптимизаций:

\paragraph {32-битные адресные значения}

Некоторые среды выполнения на 64-битных платформах используют компактные 32-битные представления ссылок для снижения потребления памяти: подобные представления используются в Android Runtime, HotSpot JVM и V8 \cite{android_art_object_reference,openjdk_compressed_oops,v8_pointer_compression}. Свободные биты указателей могут использоваться для хранения метаданных \cite{android_tagged_pointers,v8_pointer_compression}. При модификации такого компактного адреса новое значение может не помещаться в исходное представление. В подобных случаях появляются типовые компенсации: выставление нескольких старших бит адреса (часто одинаковых от обращения к обращению). Для таких случаев можно добавлять в трассу специализированные записи, которые позволят компенсировать только эти старшие биты, что может существенно улучшить репрезентативность для ряда сценариев. Подобная оптимизация требует поддержки в обработчиках трасс.

\paragraph{Хеширование адресов}

Для оптимизации компенсаций хешей адресов можно детектировать хеш функции по символьной информации (в случае полноценного вызова) и/или с помощью специализированного алгоритма, детектирующего хеширование на ГПД (при отсутствии символьной информации или инлайнинге) - с помощью методов сопоставления сигнатур графов \cite{binary_code_similarity_survey,graph_matching_survey}, символьного исполнения \cite{king_symbolic_execution} и SAT решателей \cite{handbook_of_satisfiability}.

\subsubsection{Оптимизация по памяти}

В разделе \ref{sec:real_trace_tests} приведено потребление RAM памяти во время тестовых запусков на реальных трассах исполнения заказчика (таблица \ref{tab:real_traces_tests}, также приведено усредненное значение). Подобное потребление не позволяет обрабатывать трассы с количеством инструкций порядка $1\cdot10^{9}$ и более при доступном объеме RAM, характерном для инфраструктуры заказчика. Также высокое потребление памяти существенно ограничивает возможности для параллельной обработки наборов трасс. Также в разделе \ref{sec:real_trace_tests} было показано, что основную долю потребляемой RAM занимает ГПД.

Для снижения потребления RAM памяти можно использовать специализированные библиотеки для работы с большими графами \cite{graphchi}. Такие решения хранят графы на диске и подкачивают данные в RAM при необходимости. Также могут быть применены менее специализированные библиотеки: библиотека STXXL \cite{stxxl}, предоставляющая реализации для типовых C++ контейнеров, оптимизированных под большие объемы данных; библиотеки LLNL Metall \cite{llnl_metall} и LLNL Umap \cite{llnl_umap}, предоставляющие высокоуровневые интерфейсы для выделения памяти через отображаемые файлы (ОС может выгружать такую память на диск) и контроля над выгрузкой страниц соответственно, что позволяет использовать эти библиотеки для ограничения используемой RAM.

Другим вариантом решения проблемы является отказ от хранения ГПД для всей трассы. Возможна блочная обработка трассы, со сбором ГПД только для обрабатываемого в данный момент блока. Такой подход потребует многократного сбора одних и тех же участков графа для осуществления всех необходимых обходов графа и обходов трассы, использующих данные графа. Подобное многократное восстановление данных может существенно замедлить обработку трасс, в таком случае потребуются дополнительные оптимизации.

\subsubsection{Поддержка векторных регистров}

В данный момент на ГПД добавляются только зависимости по памяти и регистрам общего назначения. При этом векторные регистры также могут использоваться для формирования адресов \cite{arm_sve_introduction}. Добавление информации о значениях векторных регистров на ГПД позволит улучшить качество генерируемых модификаций (можно будет отследить адресные значения на векторных регистрах) и поддержать инструкции, напрямую использующие эти регистры для формирования адресов (например scatter store и gather load инструкции из расширения SVE системы команд ARM64 \cite{arm_sve_introduction}).

\subsubsection{Модификация криптографических подписей адресов}

Представленное решение может быть использовано не только для изменения образа памяти трасс, но и для других задач, требующих поиска адресов в трассе. Одной из подобных задач является замена криптографических подписей адресов.

Криптографические расширения систем команд, такие как расширение PAuth архитектуры ARM64 \cite{arm_pointer_authentication,pacman}, позволяют подписывать используемые адреса и верифицировать их подпись непосредственно перед использованием. Это позволяет защитить систему от ряда кибер атак, основанных на подмене адресов потока управления.

При анализе трасс, содержащих PAuth инструкции, появляется необходимость переподписать адреса из трассы, так как ключ шифрования трассируемой системы может быть неизвестен. Локации подписанных адресов можно искать с помощью описанных в данной работе алгоритмов, после их незначительной корректировки.
